% Propietario conservado: /Users/ruben/Documents/ChatGPT/jueces y controles/REVISION_ARGUMENTAL_APERTURAS_CIERRES_20260915/IV/ES/sections/iv_apendice_enumeraciones.tex
% RESULTADO_RECUPERADO. Fuentes conservadas en la entrega:
% antecedentes/controles_articulo_I/FIRMAS_EMISION.json:
% plus_classes, times_classes, seed_id_formula y coordinate_convention.
% antecedentes/controles_articulo_I/procedencia_excepcional.md:
% sección «Comprobación finita ejecutada», matriz A y seis generadores.
% Se explicitan las enumeraciones ya utilizadas; no se cambia su dominio.
\section{Énumérations finies de l’émission et des étoiles de Witt}
\label{iv:iv:apendice-enumeraciones}

Les deux énumérations de cet appendice correspondent à des objets distincts.
La première classe les signatures des curseurs de l’émission APP ;
la seconde calcule le groupe engendré par les étoiles du plan de Witt.
Leurs domaines et leurs règles sont conservés, de sorte que chaque tableau puisse
être reproduit sans être utilisé comme entrée du calcul.

\subsection{Indices des conditions initiales et récurrence des signatures}
\label{iv:iv:apendice-firmas}

Soient \(i,j\in\{0,\ldots,8\}\) et \(a\in\{0,1,2,3\}\), avec l’ordre
des directions \((N,E,S,O)\) et les vecteurs
\[
 v_N=(-1,0),\quad v_E=(0,1),\quad
 v_S=(1,0),\quad v_O=(0,-1).
\]
L’indice
\[
 \operatorname{id}(i,j,a)=4(9i+j)+a
\]
parcourt exactement \(\{0,\ldots,323\}\). Son inverse s’obtient
au moyen de \(a=\operatorname{id}\bmod4\),
\(j=\lfloor\operatorname{id}/4\rfloor\bmod9\) et
\(i=\lfloor\operatorname{id}/36\rfloor\).
Deux indices, un par feuillet, énumèrent les \(324^2\) conditions conjointes
de \eqref{iv:eq:semillas} ; ce domaine n’est pas identifié aux 81 positions.

Pour chaque indice, le curseur est initialisé en \((i,j)\), avec le vecteur \(v_a\).
On parcourt les instants \(t=1,\ldots,54\), répartis en six fenêtres
de neuf instants. Écrivons \(r_t=1+((t-1)\bmod9)\). La règle est :
\begin{enumerate}
\item Dans le feuillet additif, on agit lorsque \(r_t\in\{1,4,7\}\) :
on ajoute \(\rho_9(i+j+2)\) à l’accumulateur de la fenêtre et on
déplace le curseur de \(v_a\), modulo neuf.
\item Dans le feuillet multiplicatif, on agit lorsque \(r_t\in\{2,5,8\}\) :
on ajoute \(\ell_9(\rho_9((i+1)(j+1)))\) et on déplace le curseur
de \(v_a\), modulo neuf. On conserve l’extension par zéro de
\(\ell_9\) hors des unités, déclarée dans la section~\ref{iv:sec:extension}.
\item Aux autres instants, ce curseur ne se déplace pas. Après
les instants \(27\) et \(54\), on remplace \(v_a\) par \(-v_a\).
\item À la fermeture de chaque fenêtre, on enregistre son accumulateur modulo dix
et on réinitialise uniquement l’accumulateur ; ni la position ni l’orientation.
\end{enumerate}
Cela définit deux fonctions sur les 324 indices,
\[
 s:\{0,\ldots,323\}\longrightarrow\{0,\ldots,9\}^6,\qquad
 e:\{0,\ldots,323\}\longrightarrow\{0,\ldots,9\}^6.
\]
La multiplicité d’une signature est le nombre d’indices ayant cette image.
L’indice témoin des tableaux est le plus petit de sa fibre.
Les préimages complètes sont récupérées en appliquant la récurrence aux
324 indices et en conservant ceux qui produisent la signature indiquée.

\begin{longtable}{rcrr}
\caption{Signatures additives : image de \(s\), multiplicité et indice témoin.}
\label{iv:iv:tabla-firmas-aditivas}\\
\toprule
Classe & Signature & Multiplicité & Indice témoin\\
\midrule
\endfirsthead
\toprule
Classe & Signature & Multiplicité & Indice témoin\\
\midrule
\endhead
1&122564&18&17\\
2&122898&18&24\\
3&212645&18&5\\
4&212988&18&0\\
5&221456&18&29\\
6&221889&18&12\\
7&456221&18&28\\
8&465898&18&21\\
9&546988&18&9\\
10&564122&18&16\\
11&645212&18&4\\
12&654889&18&33\\
13&889221&18&13\\
14&889654&18&32\\
15&898122&18&25\\
16&898465&18&20\\
17&988212&18&1\\
18&988546&18&8\\
\bottomrule
\end{longtable}

\begin{longtable}{rcrr}
\caption{Signatures multiplicatives : image de \(e\), multiplicité et indice témoin.}
\label{iv:iv:tabla-firmas-multiplicativas}\\
\toprule
Classe & Signature & Multiplicité & Indice témoin\\
\midrule
\endfirsthead
\toprule
Classe & Signature & Multiplicité & Indice témoin\\
\midrule
\endhead
1&000000&108&8\\
2&177393&8&1\\
3&177555&8&54\\
4&177771&8&11\\
5&339555&8&6\\
6&339771&8&15\\
7&339933&8&59\\
8&393177&8&0\\
9&393393&8&45\\
10&393555&8&40\\
11&555177&8&52\\
12&555339&8&4\\
13&555393&8&41\\
14&555555&24&84\\
15&555717&8&14\\
16&555771&8&18\\
17&555933&8&24\\
18&717555&8&12\\
19&717717&8&21\\
20&717933&8&19\\
21&771177&8&9\\
22&771339&8&13\\
23&771555&8&16\\
24&933339&8&57\\
25&933555&8&26\\
26&933717&8&17\\
\bottomrule
\end{longtable}

Les sommes des multiplicités sont \(18\cdot18=324\) et
\(24\cdot8+24+108=324\). Pour obtenir l’émission conjointe, on applique
\eqref{iv:eq:acoplamiento} à chaque paire de signatures. L’inverse par centaines
et dizaines démontré dans la section~\ref{iv:sec:extension} identifie les 468
émissions et donne leurs multiplicités par produit des deux fibres.
L’application ultérieure \eqref{iv:eq:psi-catalogo} produit les 243
mots visibles. Ainsi sont spécifiés séparément le domaine initial,
l’émission, ses fibres et le quotient ternaire.

\subsection{Générateurs explicites et clôture des étoiles}
\label{iv:iv:apendice-estrellas}

On conserve la matrice \(A_W\) de \eqref{iv:exc:aw} et on forme,
dans l’ordre lexicographique, les mots
\[
 c(w)=(w,wA_W),\qquad w\in\mathbb F_3^6.
\]
On les regroupe par poids et on retient les supports de poids six.
Cela produit l’énumérateur \eqref{iv:exc:enumerador} et les 132 hexades
de \(\mathcal H\). Pour chaque quadruplet \(B\subset\{1,\ldots,12\}\),
les quatre paires \(H\setminus B\), avec \(B\subset H\in\mathcal H\),
définissent \(\sigma_B\) : il fixe \(B\) et échange les deux points
de chaque paire. On obtient ainsi la famille de 495 étoiles de
la proposition~\ref{iv:exc:m12}.

Les six générateurs utilisés dans leur clôture sont exprimés par leurs listes
d’images de \(1,\ldots,12\) :
\[
\begin{array}{c|rrrrrrrrrrrr}
 &1&2&3&4&5&6&7&8&9&10&11&12\\ \hline
g_1&1&2&3&4&12&11&9&10&7&8&6&5\\
g_2&1&2&3&12&5&10&11&9&8&6&7&4\\
g_3&1&2&3&11&10&6&8&7&12&5&4&9\\
g_4&1&2&12&4&5&9&8&7&6&11&10&3\\
g_5&1&12&3&4&5&8&10&6&11&7&9&2\\
g_6&12&2&3&4&5&7&6&11&10&9&8&1
\end{array}
\]
On vérifie que chaque ligne est une étoile en la comparant aux
\(\sigma_B\) déjà construites. Pour \(j=1,\ldots,6\), la clôture
est calculée au moyen de
\[
 G_j^{(0)}=\{\operatorname{id}\},\qquad
 G_j^{(n+1)}
 =G_j^{(n)}\cup
   \{g_i\circ p:1\leq i\leq j,\ p\in G_j^{(n)}\}.
\]
La suite croissante se stabilise dans l’ensemble fini
\(\operatorname{Sym}(12)\). Comme les générateurs sont des involutions,
sa valeur stable est exactement le groupe \(G_j\) qu’ils engendrent.
Avec la convention \((g_i\circ p)(a)=g_i(p(a))\), leurs ordres sont
\[
\begin{array}{c|rrrrrr}
 j&1&2&3&4&5&6\\ \hline
 |G_j|&2&6&18&360&7920&95040 .
\end{array}
\]
On vérifie enfin que chacune des 495 étoiles appartient
à \(G_6\) et qu’elles préservent toutes \(\mathcal H\). La première inclusion,
jointe à l’appartenance des six générateurs à la famille,
établit l’égalité des deux groupes engendrés. La préservation se vérifie
au moyen de \(\{\sigma_B(H):H\in\mathcal H\}=\mathcal H\).
C’est l’énumération finie utilisée dans la preuve de
\ref{iv:exc:m12} ; la reconnaissance ultérieure comme \(M_{12}\)
conserve la référence externe qui y est déclarée.

